\documentclass[Total.tex]{subfiles}
\begin{document}
	\chapter{Simplicial Method}
		Reference:
		\begin{itemize}
			\item \url{ https://www.sas.rochester.edu/mth/sites/doug-ravenel/otherpapers/Goerss-Jardine2.pdf}
			\item \url{ https://www.sas.rochester.edu/mth/sites/doug-ravenel/otherpapers/Land-Cat.pdf}
			\item \url{https://www.math.uchicago.edu/~may/CONCISE/ConciseRevised.pdf}
		\end{itemize}
	
		\section{Abstract Simplicial Complex}
			\subsection{Simplicial Simplex}
				\begin{Def}
					\begin{itemize}
						\item 
						Given: A collection of vertices $V$. An \textbf{$n$-simplex} is a finite subset with $n+1$ different vertices:
						\begin{gather*}
							\sigma=\{v_0,\dots,v_n\}\subseteq V
						\end{gather*}
						\item Define:
						\begin{gather*}
							dim(\sigma):=n
						\end{gather*}
					\end{itemize}
				\end{Def}
				
				\subsubsection{Face}
				
				\begin{Def}
					Then, the \textbf{face} of simplex is given by the subsets of $\sigma$. 
					\begin{gather*}
						\tau\subseteq \sigma
					\end{gather*}
					
					The $i$th face is written as:
					\begin{gather*}
						d_i\sigma :=\{v_1,\dots,\hat{v_i},\dots,v_n\}
					\end{gather*}
				\end{Def}
				
				And take the following notion:
				\begin{gather*}
					X_n:=\{\sigma\in X||\sigma|=n\}
				\end{gather*}
			\subsection{Simplicial Complex}
				\begin{Def}
					The \textbf{simplicial complex} is a collection of simplicial simplices, such that:
					\begin{gather*}
						\sigma\in K\Rightarrow d_i\sigma\in K
					\end{gather*}
				\end{Def}
				\subsection{The Model in $\RR^n$}
				\begin{Def}
					Define:
					\begin{itemize}
						\item $(k+1)$ points in $\RR^n$ are in \textbf{general position} if they are not contained in a $(k-1)$-dimensional affine hyperplane in $\RR^n$;
						
						i.e. The vectors generated by difference between them form a basis of $\RR^n$;
						\item 
						\begin{gather*}
							Hull(S):=\{\sum_{x\in S} \lambda_x\cdot x|\lambda_x\in \RR_{\geq 0},\sum \lambda_x=1\}\qquad S\subseteq \RR^n,\mbox{finite}
						\end{gather*}
						\item A $k$-simplex is the convex hull of $k+1$ points that are in general position.
					\end{itemize}
				\end{Def}
				\subsection{Chain, Boundary, Cycles}
					\begin{Def}
						\begin{itemize}
							\item 
							\item 
							\item 
						\end{itemize}
					\end{Def}
				\subsection{Chain Complex by Orientation of Simplices}
					\begin{Def} Let $X$ be a simplicial complex:
						\begin{itemize}
							\item Define:
							\begin{gather*}
								C^{simp}_n(X):=\mathbb{Z}_2^{X_n}
							\end{gather*}
							\item Boundary map
							\begin{gather*}
								\partial_n:C^{simp}_n(X)\rightarrow C_{n-1}^{simp}(X)\qquad \sum a_i\sigma_i\mapsto \sum a_i\cdot \partial\sigma_i
							\end{gather*}
							\item Sequence $C^{simp}_\bullet$ is called the \textbf{simplicial chain complex} of $X$.
						\end{itemize}
					\end{Def}
					
					\begin{proposition}
						内容...
					\end{proposition}
					\begin{proof}
						内容...
					\end{proof}
				\section{Simplicial Sets}
			\begin{Def}
				Define the category $\Delta$:
				\begin{itemize}
					\item $Ob(\Delta)=\{[n]=\{0<\dots<n\}|n\in \ZZ\}$
					\item Morphisms: The order-preserving maps $[m]\rightarrow [n]$;(Not strict)
					\item Two types of fundamental morphisms:
					\begin{itemize}
						\item Coface Maps: $0\leq i\leq n$;
						\begin{gather*}
							d^i(0\rightarrow\dots\rightarrow i-1\rightarrow i\rightarrow\dots\rightarrow n-1):=(0\rightarrow \dots \rightarrow i-1\rightarrow i+1\rightarrow\dots\rightarrow n)
						\end{gather*}
						\item Codengeneracy Map: $0\leq j\leq n$
						\begin{gather*}
							s^j(0\rightarrow \dots\rightarrow j\rightarrow j+1\rightarrow\dots\rightarrow n+1):=(0\rightarrow\dots\rightarrow j\rightarrow  j\rightarrow\dots\rightarrow n)
						\end{gather*}
						
						(Remark: The notion above for the value of $d^i$, for example $j\mapsto j+1$ etc.)
					\end{itemize}
				\end{itemize}
			\end{Def}
			
			\begin{proposition}
				Identities:
				\begin{gather*}
					\begin{cases}
						d^jd^i=d^id^{j-1}\\
						s^jd^i=d^is^{j-1}\\
						s^jd^j=1=s^jd^{j+1}\\
						s^jd^i=d^{i-1}s^j\\
						s^js^i=s^is^{j+1}
					\end{cases}
				\end{gather*}
			\end{proposition}
			\begin{proof}
				内容...
			\end{proof}
			
			\begin{proposition}
				Under $(Mor(\Delta),\circ)$, then every morphism in $Mor(\Delta)$ is the composition of some $d^i_n$ and $s^i_m$.
			\end{proposition}
			\begin{proof}
				Firstly, we have decomposition
				\begin{gather*}
					f=s\circ t
				\end{gather*}
				
				where $s$ is surjective order-preserving map, $t$ is injective order-preserving map, then $s$ can be represented by composition of codegeneracy map, $t$ can be represented by composition of coface maps.
			\end{proof}
			\subsection{Simplicial Sets}
			\begin{Def}
				\begin{itemize}
					\item 
					The \textbf{simplicial set} is the functor
					\begin{gather*}
						X_\bullet:\Delta^{op}\rightarrow Set
					\end{gather*}
					
					equipped with 
					\begin{itemize}
						\item Face maps $X_\bullet(d^i)=d_i:X_{n-1}\rightarrow X_{n},0\leq i\leq n$;
						\item Degeneracy Maps $X_\bullet(s^i)=s_i:X_{n+1}\rightarrow X_{n}$;
					\end{itemize}
					\item The morphisms between is defined as following: Given $X_\bullet,Y_\bullet$, then $f_\bullet:X_\bullet\rightarrow Y_\bullet$ satisfies
					\begin{gather*}
						f_{n-1}\circ d_i^X=d_i^Y\circ f_n\\
						f_{n+1}\circ s_i^X=s_i^Y\circ f_n
					\end{gather*}
					\item Face is defined as
					\begin{gather*}
						f_iX_n:=im(d_i:X_{n-1}\rightarrow X_n)
					\end{gather*}
				\end{itemize}
			\end{Def}
			
			\begin{proposition}
				Identities
				\begin{gather*}
					\begin{cases}
						d_id_j=d_{j-1}d_i&i<j\\
						d_is_j=s_{j-1}d_i&i<j\\
						d_js_j=1=d_{j+1}s_j\\
						d_js_j=s_jd_{i-1}&i>j+1\\
						s_is_j=s_{j+1}s_i&i\leq j
					\end{cases}
				\end{gather*}
			\end{proposition}
			\begin{proof}
				内容...
			\end{proof}
			
			\subsubsection{Moore Complex}
				\begin{lemma}
					We have $\partial^2=0$. 
				\end{lemma}
				\begin{proof}
					内容...
				\end{proof}
				
				\begin{Def}
					Define:
					\begin{gather*}
						\partial:=\sum (-1)^k d_k
					\end{gather*}
					
					extends it to free abelian group. Then we get a chain complex
					\begin{gather*}
						\dots\rightarrow \ZZ X_{n+1}\stackrel{\partial}{\rightarrow}\ZZ X_n\stackrel{\partial}{\rightarrow} \ZZ X_{n-1}\rightarrow\dots
					\end{gather*}
				\end{Def}
				
				\subsubsection{Normalized Chain Complex}
				
				\subsubsection{Limit}
				
				\begin{proposition}
					Let $X\in sSet$. There is an isomorphism
					\begin{gather*}
						X\cong colim_{\Delta^n\rightarrow X} \Delta^n
					\end{gather*}
				\end{proposition}
				\begin{proof}
					内容...
				\end{proof}
				
				\subsubsection{Self-Enrichment}
				
				\begin{proposition}
					Category $sSet$ is self-enriched.
				\end{proposition}
				\begin{proof}
					$sSet$ is \underline{Cartesian closed}: With Cartesian product
					\begin{gather*}
						(sSet,\times,\Delta[0])
					\end{gather*}
				\end{proof}
				
			\subsection{Nerve}
				
				\subsubsection{Skeleta}
				\subsubsection{Nerve}
				\begin{Def}
					Consider a small category $\mathcal{C}$. And take the elements $[n]\in \Delta^{op}$ as a category. Then, define:
					\begin{gather*}
						(N(C):\Delta^{op}\rightarrow Set)\in sSet\qquad N(C)_n=Func([n],\mathcal{C})=Mor_{Cat}([n],C)
					\end{gather*}
					
					\textbf{Remark} Sometimes the nerve is called \textbf{classifying space}
				\end{Def}
				
				\begin{proposition}
					The functor of nerve
					\begin{gather*}
						N:Cat\rightarrow sSet
					\end{gather*}
					
					is fully faithful.
				\end{proposition}
				\begin{proof}
					内容...
				\end{proof}
				
				\begin{example}
					(Algebraic K-Space) \url{https://www.sas.rochester.edu/mth/sites/doug-ravenel/otherpapers/Goerss-Jardine2.pdf}, example 1.6, we can get the 'algebraic $K$-space' by taking the geometric relization of the classifying space $BQ\mathcal{A}$, where $\mathcal{A}$ is exact.
				\end{example}
				\subsubsection{Generalized Nerve}
				
				\subsection{Relization }
				
				\subsubsection{Nerve-Relization Adjunction}
				
				\begin{theorem}
					内容...
				\end{theorem}
				\begin{proof}
					内容...
				\end{proof}
				
				
				Actually, the relization operation is the  \textbf{left Kan extension} of $S_\mathcal{C}:S\rightarrow \mathcal{C}$ along the Yoneda embedding $S\hookrightarrow [S^{op},V]$(i.e. the Yoneda extension)
			\subsection{Geometric Relization}
				\begin{Def}
					Define: Given a simplicial set $X_\bullet$,
					\begin{gather*}
						|X_\bullet|:=(\sqcup_{n\geq 0} X_n\times \Delta^n)/\sim
					\end{gather*}
					
					where the relation is given by:
					\begin{gather*}
						()\\
						()
					\end{gather*}
					
					And: $\gamma:\Gamma X\rightarrow X$ by giving
					\begin{gather*}
						内容...
					\end{gather*}
				\end{Def}
				
				\begin{theorem}
					For any space $X$, $\Gamma X$ is a CW complex with one $n$-cell for each nondegenerate singular $n$-simplex, and the cellular chain complex
					\begin{gather*}
						C_*(\Gamma X)\cong Sing(X)_* \in Ch({\mathcal Ab})
					\end{gather*}
				\end{theorem}
				\begin{proof}
					内容...
				\end{proof}
				
				\begin{theorem}
					For any space $X$, the map $\gamma:\Gamma X\rightarrow X$ is a weak equivalence.
				\end{theorem}
				\begin{proof}
					内容...
				\end{proof}
				
				\subsubsection{Colimit}
				
				Especially, 
				
				\begin{proposition}
					The realization functor preserves colimit.
				\end{proposition}
				\begin{proof}
					The main reason is, that there exists adjunction
					\begin{gather*}
						Hom_{Top}(|X|,Y)\cong Hom_{sSet}(X,Sing(Y))
					\end{gather*}
				\end{proof}
				
				Thus, we have
				\begin{gather*}
					内容...
				\end{gather*}
			\subsection{Standard Simplicial Simplex}
				\begin{Def}
					Define: $\Delta[n]:=Hom_\Delta(-,[n])\in sSet$, i.e. 
					\begin{gather*}
						\Delta[n]_m=\Delta[n](m):=Hom_\Delta([m],[n])\in Set\qquad m\in Ob(\Delta )
					\end{gather*}
				\end{Def}
				
				\subsubsection{Geometric Realization}
				
				\begin{proposition}
					Its geometric relization is $\Delta^n$. 
					\begin{gather*}
						|\Delta[n]|\cong \Delta^n
					\end{gather*}
				\end{proposition}
				\begin{proof}
					内容...
				\end{proof}
				
				\subsubsection{Presheaf}
				
				\begin{proposition}
					(Yoneda) Suppose $X$ be a simplicial set, then
					\begin{gather*}
						X_n\cong Hom_{sSet}(\Delta[n],X)
					\end{gather*}
				\end{proposition}
				\begin{proof}
					Just apply Yoneda's lemma.(Presheaf)
					
					This bijection is specified as following: 
					
					Given some $\varphi\in Hom_{sSet}(\Delta[n],X)$. At $n$, there exists $\varphi_n$ with 
					\begin{gather*}
						\varphi_n:\Delta[n](n)=Hom_\Delta([n],[n])\rightarrow X_n
					\end{gather*}
					
					We could select the point $x:=\varphi(\iota_n)=\varphi_n(\iota_n)$
					
					Conversely, given $x\in X_n$, the morphism $\varphi^x\in Hom_{sSet}(\Delta[n],X)$ is specified by
					\begin{gather*}
						\varphi^x_n:\iota_n\mapsto x\\
						Hom_\Delta([m],[n])\ni \alpha\mapsto Y(\alpha)(x)
					\end{gather*}
				\end{proof}
			
			\subsubsection{Boundary, Horn}
			\begin{Def}
				Define:(Standard)
				\begin{itemize}
					\item Boundary $\partial\Delta[n]:=\partial(\Delta[n])$
					
					\item Horn $\Delta[n]\supseteq \Lambda^k[n]:=\cup_{0\leq i\leq n,i\not=k} f_i\Delta[n]=\cup_{i\not=k} im(\Delta[n-1]\stackrel{d_i^*}{\rightarrow} \Delta[n])$
				\end{itemize}
			\end{Def}
			
			Furthermore, given a simplicial set $X$, and given
			\begin{Def}
				The \textbf{horn} could be described as:
				\begin{gather*}
					Horn^k_n(X):=\{(x_i)_{i\not=k}|x_i\in X_{n-1}, d_ix_j=d_{j-1}x_i,i<j,i,j\not=k\}
				\end{gather*}
			\end{Def}
			
			Then
			\begin{proposition}
				\begin{gather*}
					Horn^k_n(X)=Hom_{sSets}(\Lambda^k[n],X)
				\end{gather*}
			\end{proposition}
			\begin{proof}
				Define:
				\begin{gather*}
					\iota_r:\Delta[n-1]\rightarrow \Lambda^k[n]\subseteq \Delta[n]\qquad r\not=k
				\end{gather*}
				
				which is the inclusion of $r$-face. Suppose $f\in Hom_{sSets}(\Lambda^k[n],X)$, then
				\begin{gather*}
					f\circ \iota_r\in Hom_{sSet}(\Delta[n-1],X)\cong X_{n-1}\ni x_r
				\end{gather*}
				
				Thus, the simplicial map of $d_i x_j$ is
				\begin{center}
					\begin{tikzcd}
						\Delta[n-2]\arrow[r,"d_i"]&\Delta[n-1]\arrow[r,"\iota_j"]&\Lambda^k[n]\arrow[r,"f"]&X
					\end{tikzcd},\begin{tikzcd}
						X_{n-1}\arrow[r,"d_i"]\arrow[d,"\cong"]&X_{n-2}\arrow[d,"\cong"]\\
						Hom(\Delta[n-1],X)\arrow[r]&Hom(\Delta[n-2],X)
					\end{tikzcd}
				\end{center}
				
				Similarly $d_{j-1}x_i\leftrightarrow f\circ \iota_i\circ d_{j-1}$, then
				\begin{gather*}
					\iota_i\circ d_{j-1}=\iota_j\circ d_{i}
				\end{gather*}
				
				because $\iota_i:\Delta[n-1]\rightarrow \Delta[n]$ act as $d_i$, and $d_i\circ d_{j-1}=d_j\circ d_i$. Thus $d_ix_j=d_{j-1}x_i$
			\end{proof}
			
			The graph can be described as
			\begin{center}
				\begin{tikzcd}
					&0\arrow[ld]\arrow[rd]\\
					1&&2
				\end{tikzcd}$\subset$ \begin{tikzcd}
					&0\arrow[ld]\arrow[rd]\\
					1\arrow[rr]&&2
				\end{tikzcd}$=\Delta[2]$. 
			\end{center}
			
			\textbf{Remark} $\Lambda^k[n]\subseteq \Delta[n]\in sSet$, but $Horn^k_n(X)$ is not a simplicial set. 
			
			\begin{Def}
				\begin{itemize}
					\item Inner horns: $\Lambda^k(n),0<k<n$;
					\item Outer horns: $\Lambda^0[n],\Lambda^[n]$;
				\end{itemize}
			\end{Def}
			
			\subsubsection{Horn Filler}
			
			\begin{Def}
				A \textbf{horn filler} is an extension of horn:
				\begin{center}
					
					\begin{tikzcd} \Lambda^k[n]\arrow[r,"h"]\arrow[d,hook,"\iota"'] & X\\ \Delta[n]\arrow[ru,dashed,"\widetilde h"'] \end{tikzcd}\qquad \begin{tikzcd}
						Horn^k_n(X)\\
						Hom(\Delta[n],X)\cong X_n\arrow[u,mapsto,"\rho^k_n"]
					\end{tikzcd}
				\end{center}
				
				Equivalently, a correspondence $x\in X_n$ such that $d_ix=x_i,i\not=k$. 
				
				The collection of fillers can actually be given by the fiber of $\rho^k_n$.
			\end{Def}
			
			\textbf{Remark} Filler is not unique.
			
		\section{Geometric Simplicial Complex}
		\subsubsection{Standard Simplicial Simplex}
		\begin{Def}
			\begin{itemize}
				\item 
				The $n$-dimensional standard simplex:
				\begin{gather*}
					\Delta^n=\Delta[n]=\{(t_0,\dots,t_{n})\in \RR^{n+1}|\sum t_i=1\}\subseteq \RR^{n+1}
				\end{gather*}
				\item For (weak) increasing map $[\alpha]:[m]\rightarrow [n]$ define
				\begin{gather*}
					\Delta[\alpha]:\Delta[m]\rightarrow \Delta[n]\qquad \sum t_i e_i\mapsto \sum t_i e_{\alpha(i)}
				\end{gather*}
				\item Define the maps
				\begin{gather*}
					\delta^{n}_i:[n-1]\rightarrow [n]\qquad \theta^n_i:[n]\rightarrow [n+1]
				\end{gather*}
				
				\item $d^n_i:=\Delta(\delta^n_i),s^n_i:=\Delta(\theta^n_i)$;
			\end{itemize}
		\end{Def}
		
		\begin{proposition}
			We have the following identities:
			\begin{gather*}
				内容...
			\end{gather*}
		\end{proposition}
		\begin{proof}
			内容...
		\end{proof}
		\subsection{Singular Simplicial Set}
			\begin{Def}
				Given a topological space $X$. Define:
				\begin{gather*}
					Sing(Y)_n:=Hom_{Top}(\Delta^n,Y)
				\end{gather*}
			\end{Def}
			
			\begin{proposition}
				$Sing(Y)$ is a simplicial set. With the structure:
				\begin{gather*}
					Sing(Y)_{(-)}:\Delta^{op}\rightarrow Set
				\end{gather*}
				
				Notice: $Top$ is small category.
			\end{proposition}
			\begin{proof}
				内容...
			\end{proof}
			
			\begin{corollary}
				$Sing(-)$ is a functor $Top\rightarrow sSet$.
			\end{corollary}
			\begin{proof}
				Definition.
			\end{proof}
			
			\subsubsection{Relization-$Sing$ Adjunction}
				\begin{theorem}
					There exists adjunction  $|-|\dashv Sing$,
					\begin{center}
						\begin{tikzcd}
							sSet\arrow[rr,"|-|",shift left=4]&{\small \perp}&Top\arrow[ll,"Sing",shift left=4]
						\end{tikzcd}
					\end{center}
					
					i.e. $Hom_{Top}(|X|,Y)\cong Hom_{sSet}(X,Sing(Y))$.
				\end{theorem}
				\begin{proof}
					内容...
				\end{proof}
			\subsection{Simgular Chain Complex}
				\begin{Def}
					Define:
					\begin{itemize}
						\item $d_i:$
						\item $s_i:$
					\end{itemize}
				\end{Def}
				
				\begin{proposition}
					内容...
				\end{proposition}
				
				\begin{proof}
					内容...
				\end{proof}
				
				Especially,
				\begin{gather*}
					C_*(X;\pi)L=C_*(X)\oplus \pi
				\end{gather*}
		\section{Kan Complexes}
			\subsection{Kan Fibration}
				\begin{Def}
					A \textbf{Kan fibration} is a morphism $Y\rightarrow X$ of the simplicial set with the \textbf{lifting-property} for all horn-inclusions.
					
					\begin{center}
						\begin{tikzcd}
							\Lambda^k[n]\arrow[d,"\subseteq"]\arrow[r]&Y\arrow[d,"\pi"]\\
							\Delta[n]\arrow[r]\arrow[ru,dashrightarrow]&X
						\end{tikzcd}
					\end{center}
					
					Equivalently, by powering, the  morphism
					\begin{gather*}
						Y^{\Delta[n]}\rightarrow Y^{\Lambda^k[n]}\times_{X^{\Lambda^k[n]}} X^{\Delta^k[n]}
					\end{gather*}
					
					is an epimorphism.
				\end{Def}
				
				The right-side is the fibre product
				\begin{tikzcd}
					Hom(\Delta[n],Y)\arrow[rd,dashrightarrow,"epi"]	\arrow[rdd]\arrow[rrd] \\
				&RHS	\arrow[d]\arrow[r]&Hom(\Lambda^k[n], Y)\arrow[d,"\pi^*" ]]\\
					&Hom(\Delta[n],X)\arrow[r]&Hom(\Lambda^k[n],X)
				\end{tikzcd}
				
				\begin{example}
					\begin{itemize}
						\item $n=1$:
						\item $n=2$:
					\end{itemize}
				\end{example}
				
				This will be proved in fibration, homotopy theory?
				
				Furthermore, some of definitions:
				\begin{Def}
					\begin{itemize}
						\item Acyclic Kan fibration: A Kan fibration is acyclic if it is a weak equivalence;
						\item Minimal Kan fibration: The kan fibration $p:E\rightarrow B$ is minimal, if for all cells $x,y\in \Delta[n]\rightarrow E$, the condition
						\begin{gather*}
							\partial_i x=\partial_i y,p(x)=p(y)\Rightarrow \partial_k x=\partial_k y,\forall k
						\end{gather*}
						\item Quasi-fibration: The lifting property only impose in inner horns: Such $k=0,n$.
						\item Left, Right Kan Fibration:
						\begin{itemize}
							\item Left: Lifting: Except the last one;
							\item Right lifting: Except the first one.
						\end{itemize}
					\end{itemize}
				\end{Def}
				
				\subsubsection{Surjective Kan Fibrations}
				
				\begin{lemma}
					(Criterion for Epimoprhism)As soon as a Kan fibration is a surjection in degree $0$(i.e. $\pi_0$ is surjective) then it is a surjection is all degrees. Hence $\pi$ is an epimorphism of simplicial sets.
					\begin{center}
						$f:X\rightarrow Y$ is Kan fibration, $f_0$ is surjective$\Rightarrow \forall_n f_n$ is surjective. 
					\end{center}
				\end{lemma}
				\begin{proof}
					Given $n\in \NN$, then there exists right-lifting
					\begin{center}
						\begin{tikzcd}
							&Y\\
							\Delta[n]&X
						\end{tikzcd}
					\end{center}
				\end{proof}
				 
				\textbf{Remark} A short proof is by \textbf{the conclusion of model category}. 
				
				Notice: 
				
				\begin{lemma}
					When: $f$ is a Kan fibration, which is surjective on the connected components $\pi_0(f)$, then it is a surjection is  a surjection in all degree.
				\end{lemma}
				\begin{proof}
					内容...
				\end{proof}
				
				\begin{proposition}
					Let $f:X\rightarrow Y$ be any morphism of simplicial sets which is a surjection on connected components
					\begin{gather*}
						\pi_0(f):\pi_0(X)\xrightarrow \pi_0(Y)
					\end{gather*}
					
					Then, $f$ factors as a simplicial weak equivalence followed by a Kan fibratiuon which is a degreewise surjection.
				\end{proposition}
				\begin{proof}
					内容...
				\end{proof}
				\subsubsection{Acyclic Kan fibration}
				\subsubsection{Pullback,Homotopy Pullback}
			\subsection{Kan Complexes}
				\begin{lemma}
					(\textbf{Kan Condition})The followings are equivalent:
					\begin{itemize}
						\item All horns of $S$ have fillers;
						\item Equivalently, the unique homomorphism $S\rightarrow pt$ from $S$ to the point space $\bullet$ (the terminal simplicial set) is a Kan fibration;
						\item Equivalently, for all diagrams of the form, there exists extension
						\begin{center}
							\begin{tikzcd}
								{\Lambda^i[n]}\arrow[r]\arrow[d,hook]&S\\
								{\Delta[n]}
							\end{tikzcd}$\Rightarrow$ \begin{tikzcd}
								{\Lambda^i[n]}\arrow[r]\arrow[d,hookrightarrow]&S\\
								{\Delta[n]}\arrow[ru]
							\end{tikzcd}\\
							Equivalently \qquad \begin{tikzcd}
									{\Lambda^i[n]}\arrow[r]\arrow[d]&S\arrow[d]\\
								{\Delta[n]}\arrow[r]&\bullet
							\end{tikzcd}$\Rightarrow$ \begin{tikzcd}
								{\Lambda^i[n]}\arrow[r]\arrow[d]&S\arrow[d]\\
							{\Delta[n]}\arrow[r]\arrow[ru,dashrightarrow]&\bullet
							\end{tikzcd}
						\end{center}
						\item The map from $n$-simplicies to $(n,i)$ horn
						\begin{gather*}
							[\Delta[n],S]\twoheadrightarrow [\Lambda^i[n],S]
						\end{gather*}
						
						are surjective submersions.
					\end{itemize}
				\end{lemma}
				\begin{proof}
					\begin{itemize}
						\item (1)$\Rightarrow$ (2)
						\item (2)$\Rightarrow$ (3)
						\item (3)$\Rightarrow$ (4)
					\end{itemize}
				\end{proof}
				
				\begin{lemma}
					(Weak Kan Condition) The followingss are equivalent:
					\begin{itemize}
						\item All inner horns in $C$ have fillers. 
						\item The morphism of simplicial sets
						\begin{gather*}
							sSet(\Delta[2],C)\rightarrow sSet(\Lambda^1[2],C)
						\end{gather*}
						
						is an acyclic Kan fibration.
					\end{itemize}
				\end{lemma}
				\begin{proof}
					\begin{itemize}
						\item (1)$\Rightarrow$ (2)
						\item (2)$\Rightarrow$ (1)
					\end{itemize}
				\end{proof}
				
				
				\begin{Def}
					
					\begin{itemize}
						\item Kan Commplex: A \textbf{Kan complex} is a simplicial set $S$ that satisfies the Kan condition
						\item Weak Kan Complex	/Quasi-Category: A simplicial set $S$ satisfies the weak Kan condition.
						
						Especially, if the quasi-category equipped with a choice of inner horn fillers.
					\end{itemize}
				\end{Def}
				
			\subsection{Presentation of $\partial \Delta[n],\Lambda^n_k$}
				The presentation:
				\begin{center}
					\begin{tikzcd}
						内容...
					\end{tikzcd}\begin{tikzcd}
					内容...
					\end{tikzcd}
 				\end{center}
			\subsection{Conclusions}
			\subsubsection{Every Topological Space is Kan Complex}
				\begin{proposition}
					内容...
				\end{proposition}
				\begin{proof}
					内容...
				\end{proof}
			\subsubsection{Nerve is Kan Complex$\Leftrightarrow$ Groupoid}
				\begin{proposition}
					The followings are equivalent:
					\begin{itemize}
						\item $N(\mathcal{C})$ is a chain complex;
						\item $\mathcal{C}$ is a groupoid.
					\end{itemize}
				\end{proposition}
				\begin{proof}
					内容...
				\end{proof}
		\section{Anodyne Extension}
		\section{Simplicial Homotopy}
		Recall
		\begin{Def}
			\begin{itemize}
				\item (Homotopy) Given simplicial maps $f,g:K\rightarrow X$, the a \textbf{simplicial homotopy} $f$ to $g$ is the commutative diagram
				\begin{center}
					\begin{tikzcd}
						内容...
					\end{tikzcd}
				\end{center}
				\item (Relative Homotopy)
			\end{itemize}
		\end{Def}
		
		\begin{lemma}
			内容...
		\end{lemma}
		\begin{proof}
			内容...
		\end{proof}
		\section{Dold-Kan Correspondence}
			\subsection{Abelian Case}
				\begin{theorem}
					(Dold-Puppe) Suppose $A$ be an abelian category. There exists an euqivalence of categories
					\begin{center}
						\begin{tikzcd}
							A^{\Delta_{op}}&Ch^+_\bullet(A)
						\end{tikzcd}
					\end{center}
				\end{theorem}
				\begin{proof}
					内容...
				\end{proof}
				
				\begin{theorem}
					(Kan)
				\end{theorem}
				\begin{proof}
					内容...
				\end{proof}
		\section{Simplicial Resolutions}
	\chapter{Homology}
		
		
		\section{Axioms for Homology}
		\subsection{Cellular Homology}
		
		Firstly, we have: The homeomorphism
		\begin{gather*}
			v_n:D^n/S^{n-1}\rightarrow S^n\\
			\iota_n:S^n\rightarrow \Sigma S^{n-1}
		\end{gather*}
		
		Define the homeomorphism of pairs
		\begin{gather*}
			内容...
		\end{gather*}
		
		\begin{lemma}
			内容...
		\end{lemma}
		\begin{proof}
			内容...
		\end{proof}
		\subsection{Axioms}
		\begin{theorem}
			For any integer $q$, there exists functors $H_q(X,A;\pi)$ from the homotopy category of pairs of spaces $Ho((X,A))$ to the category of abelian groups, together with natural transformations
			\begin{gather*}
				\partial:H_q(X,A;\pi)\rightarrow H_{q-1}(A,\pi)
			\end{gather*}
			
			where $H_q(X,\pi):=H_q(X,\emptyset;\pi)$. These satsifies the following \underline{axioms}
			
			\begin{itemize}
				\item (Dimension) If $X=\bullet$, then
				\begin{gather*}
					H_0(\bullet;\pi)=\pi\qquad H_q(\bullet;\pi)=0,q\geq 1
				\end{gather*}
				\item (Exactness) The following sequence is exact, for inclusion
				\begin{gather*}
					A\hookrightarrow X\qquad (X,\emptyset)\rightarrow (X,A)
				\end{gather*}
				
				we have:
				\begin{gather*}
					\dots\rightarrow H_n(X,A;\pi)\rightarrow H_n(X;\pi)\rightarrow H_n(X,A,\pi)\stackrel{\partial}{\rightarrow} H_{q-1}(A,\pi)\rightarrow\dots
 				\end{gather*}
				
				called the \textbf{connecting homomorphism}
				\item (Excision) If $(X;A,B)$ is an excisive triad, so that $X$ is the union of the iterior of $A$ and $B$,
				\begin{gather*}
					X=\dot{A}\cup \dot{B}
				\end{gather*}
				
				then the inclusion $(A,A\cap B)\rightarrow (X,B)$ induces an isomorphism:
				\begin{gather*}
					H_*(A,A\cap B;\pi)\rightarrow H_*(X,B;\pi)
				\end{gather*}
				\item (Additivity) If $(X,A)$ is a disjoint union of a set of pairs $(X_i,A_i)$, then
				\begin{gather*}
					\sum_i H_*(X_i,A_i;\pi)\rightarrow H_*(X,A,\pi)
				\end{gather*}
				
				is an isomorphism induced by $(X_i,A_i)\hookrightarrow H_*(X,A;\pi)$
				\item (Weak Equivalence) If $f:(X,A)\rightarrow (Y,B)$ is a weak equivalence, then
				\begin{gather*}
					f_*:H_*(X,A,\pi)\rightarrow H_*(Y,B;\pi)
				\end{gather*}
				
				is an isomorphism.
			\end{itemize}
		\end{theorem}
		\begin{proof}
			We just use the following conclusion on CW-complex and CW-approxiamation theorem, and the conclusions are only related to the values of homology.
			
			Weak equivalence:
			
			Additivity:
			
			Excision:
		\end{proof}
		
		\begin{theorem}
			(CW-version) For integer $q$, there exist functor $H_q(X,A;\pi):Ho(CW-Complexes)\rightarrow Ab$ together with a natural transformation
			\begin{gather*}
				\partial:H_q(X,A)\rightarrow H_{q-1}(A,\pi)
			\end{gather*}
			
			And axioms:
			\begin{itemize}
				\item Dimension: If $X$ is a point, then
				\begin{gather*}
					内容...
				\end{gather*}
				\item Exactness: The following sequence is exact, induced by $A\rightarrow X,(X,\emptyset)\rightarrow (X,A)$:
				
				\begin{gather*}
					\dots\rightarrow 
				\end{gather*}
			
				\item Excision: If $X$ is the union of subcomplexes $A,B$, then the inclusion $(A,A\cap B)\rightarrow (X,B)$ induces an isomorphism
				\begin{gather*}
					H_*(A,A\cap B;\pi)\rightarrow H_*(X,B;\pi)
				\end{gather*}
				\item Additivity: If $(X,A)$ is the disjoint union of a set of pairs $(X_i,A_i)$, then the inclussions $(X_i,A_i)\rightarrow (X,A)$ induce an isomorphism
				\begin{gather*}
					内容...
				\end{gather*}
			\end{itemize}
		\end{theorem}
		\begin{proof}
			Direct from the discussion of CW-complex.
		\end{proof}
		\section{Reduced Homology}
			\begin{Def}
				Suppose $(X,*)$ be a pointed space, then
				\begin{gather*}
					\tilde{E}_q(X):=E_q(X,*)
				\end{gather*}
			\end{Def}
			
			Use Mayer-Vietoris by selecting
			\begin{gather*}
				内容...
			\end{gather*}
			\subsection{Quotient}
			
			
			What's difference? We have:
			\begin{proposition}
				内容...
			\end{proposition}
			\begin{proof}
				内容...
			\end{proof}
			\subsection{Cofibration, Homology Pairs}
			\begin{theorem}
				For any cofibration $i:A\rightarrow X$, the quotient map $q:(X,A)\rightarrow (X/A,*)$ induces an isomorphism
				\begin{gather*}
					内容...
				\end{gather*}
			\end{theorem}
			\begin{proof}
				内容...
			\end{proof}
			\subsection{Suspension, Long Exact Sequence}
			
			\begin{theorem}
				For a \textbf{nondegenerately base space} $X$, there exists a natural issomorphism
				\begin{gather*}
					内容...
				\end{gather*}
			\end{theorem}
			\begin{proof}
				内容...
			\end{proof}
			
			\begin{corollary}
				Let $*\in A\subseteq X$, then the inclusion of pointed spcaes $i:A\rightarrow X$ is a cofibration between nondegenerately base spaces. In teh long exact sequence
				\begin{gather*}
					内容...
				\end{gather*}
			\end{corollary}
			\begin{proof}
				内容...
			\end{proof}
			
			\begin{corollary}
				For any $n,q$ we have
				\begin{gather*}
					\tilde{S^n}\cong \tilde{E}^{q-n}(*)
				\end{gather*}
			\end{corollary}
			\begin{proof}
				内容...
			\end{proof}
			\subsection{Axioms}
			\begin{Def}
				内容...
			\end{Def}
			
			\begin{Def}
				(CW)
			\end{Def}
			\subsection{Mayer-Vietoris}
			
			\begin{proposition}
				内容...
			\end{proposition}
			\begin{proof}
				内容...
			\end{proof}
			
			\begin{lemma}
				内容...
			\end{lemma}
			\begin{proof}
				内容...
			\end{proof}
			
			\begin{theorem}
				(Relative Mayer-Vietoris Sequence)
			\end{theorem}
		\section{Simplicial Sets Homology}
		\section{Singular Homology}	\subsection{Singular Homology Groups}
			\begin{Def}
				内容...
			\end{Def}
			
			\subsubsection{Relative Singular Homology Groups}
			\subsubsection{Singular Homology-Simplicail Sets Homology}
				\begin{theorem}
					There exists an isomorphism
					\begin{gather*}
						内容...
					\end{gather*}
				\end{theorem}
				\begin{proof}
					内容...
				\end{proof}
			\subsection{Hurewicz's Theorem in Dimension $0$}
			
			\begin{theorem}
				内容...
			\end{theorem}
			\begin{proof}
				内容...
			\end{proof}
			\subsection{Hurewicz's Theorem in Dimension $1$}
				\begin{theorem}
					内容...
				\end{theorem}
				\begin{proof}
					内容...
				\end{proof}
			\subsection{Eilenberg-Zilber Theorem}
			\subsection{Weak Equivalences}
				\subsubsection{Eilenberg Subcomplex}
					\begin{Def}
						Define:
						\begin{gather*}
							内容...
						\end{gather*}
					\end{Def}
		
		\subsection{Suspension on Homology}
		\section{Local Homology}
		\section{Product}
		\section{Hurewicz's Theorem}
			\begin{Def}
				(Hurewicz Homomorphism) Suppose $(X,x)$ be pointed space, then the \textbf{Hurewicz homomorphism} is the function
				\begin{gather*}
					\Phi:\pi_k(X,x)\rightarrow H_k(X)\\
					[S^k,X]\ni f\mapsto f_*[S_k]
				\end{gather*}
				
				Notice that, $\Phi$ is a natural transformation
				\begin{center}
					\begin{tikzcd}
						Top_\bullet\arrow[rr,bend left=20]\arrow[rr,bend right=20]&\Downarrow&{\mathcal Ab}
					\end{tikzcd}
				\end{center}
			\end{Def}
			
			\subsection{Hurewicz Theorem}
				\begin{theorem}
					(Hurewicz) 
				\end{theorem}
				\begin{proof}
					内容...
				\end{proof}
			\subsection{Theorem for Cohomology by UCT}
	\chapter{Cohomology}
		\section{Axiom for Cohomology}
		\section{Cup and Cap Product}
		\section{Cohomology Relations}
		\section{Cohomology Operators}
			\subsection{Cohomology Operations}
			\begin{Def}
				Let: $A,B$ be abelian groups, and $n,m$ natural numbers, then a \textbf{cohomology operation} of \underline{type} $(A,n,B,m)$ is 
				\begin{gather*}
					\tau_{(-)}=\tau:H^n(-,A)\rightarrow H^m(-,B)
				\end{gather*}
				
				Denoted by $Oper(A,n,B,m)$.
				
				\textbf{Remark} $\tau_X$ does not demand, that it is a group homomorphism. But, we can take the pointwise addtion of the functions.
			\end{Def}
			
			\begin{lemma}
				We have
				\begin{itemize}
					\item $\sigma,\tau\in Oper(A,n,B,m)\Rightarrow \sigma+\tau\in Oper(A,n,B,m)$;
					\item Thus, $Oper(A,n,B,m)$ is an abelian group.
				\end{itemize}
			\end{lemma}
			\begin{proof}
				内容...
			\end{proof}
			
			\begin{Def}
				Firstly, we have Hurewicz isomorphism $\pi_n(K(A,n),*)\cong A$. Then, the \textbf{fundamental class} $\iota_{n,A}\in H_n(K(A,n),A)$ is the unique element such that the composition
				\begin{gather*}
					内容...
				\end{gather*}
			\end{Def}
			
			\begin{theorem}
				There exists an isomorphism
				\begin{gather*}
					Oper(A,n,B,m)\rightarrow H^m(K(A,n),B)\qquad \tau\mapsto \tau(\iota_{n,A})
				\end{gather*}
			\end{theorem}
			
			\subsubsection{Yoneda Representability}
			
			Recall
			\begin{itemize}
				\item CW Approxiamation:
				\item 
			\end{itemize}
			
			\subsubsection{$\cup_1$-Product}
			\begin{Def}
				内容...
			\end{Def}
			\subsection{Stable Cohomology Operations}
			\subsubsection{Reduced Cohomology}
			\begin{Def}
				内容...
			\end{Def}
			\subsection{Cohomology in Stable Range}
			\subsection{Adem Relation }
	\chapter{Homology and Cohomology}
		\section{Euler Characteristics}
			\subsection{CW Euler-Characteristics}
		\section{Poincare Duality Theorem}
			Recall:
			\begin{itemize}
				\item Natural Homomorphism of Dual
				\begin{gather*}
					内容...
				\end{gather*}
				\item 
				\item Cup Product:
				\begin{gather*}
					内容...
				\end{gather*}
			\end{itemize}
			
			\begin{theorem}
				(Poincare Duality Theorem)
			\end{theorem}
			
		\section{Orientation, Fundamental Classes}
		\section{Orientation Cover}
\end{document}